\subsection{Верификация модели}

Проверка выполнена 28 октября 2026 года в Windows с использованием
SPIN 6.5.1 и GCC 16.2.0. Для запуска использовался рассмотренный скрипт
\code{run.py}. В отдельной папке запуска были сохранены копии исходных
файлов, сгенерированный верификатор и логи всех этапов.

Сначала выполнялась структурная проверка без \code{never claim},
затем --- отдельный запуск для каждой из 14 формул. Такой порядок
позволяет разделить недопустимое конечное состояние самой модели
и нарушение конкретного требования. При проверке использовались полный
поиск, сжатие \code{COLLAPSE} и стандартное сокращение частичных порядков.

В отличие от запуска одного сценария, верификация исследует допустимые
варианты запросов и взаимодействия процессов. Память не сбрасывается
между транзакциями, поэтому проверяются также последующие обращения
к уже измененным регистрам. Начальное состояние при этом фиксировано:
все регистры равны нулю, Master находится в состоянии подготовки.

\subsubsection{Результаты автоматического и отдельного запусков}

Автоматический запуск выполнил 15 проверок и сформировал итог:
\code{PASS=14, FAIL=0, ERROR=1}. Структурная проверка и 13 LTL-свойств
завершились без нарушений. Для свойства \code{p12_transaction_progress}
сработало ограничение Python-скрипта в 180 секунд: в этом запуске SPIN
выполнил около 35 миллионов посещений состояний, но обход не был закончен.
Поэтому результат \code{ERROR} отражал недостаточное время проверки,
а не найденный контрпример.

Чтобы получить законченный результат для свойства прогресса, тот же
собранный верификатор был запущен отдельно из папки проверки. Параметры
поиска, модель и LTL-формула сохранены; ограничение времени, установленное
в Python-обертке, к этой команде не применялось.

\listing[]{manual_command.txt}{Отдельная проверка свойства прогресса}

Отдельный поиск завершился за 214 секунд с результатом \code{errors: 0}.
Было сохранено 28\,849\,450 состояний и выполнено около 43,1 миллиона
их посещений. Таким образом, все 14 LTL-свойств получили законченный
результат без нарушений; структурная проверка также пройдена. Контрпримеров
не найдено. Итоги завершенных проверок приведены в таблице~\ref{tab:verification}.

\begin{table}[H]
\centering
\caption{Результаты завершенных проверок модели}
\label{tab:verification}
{\small\setlength{\tabcolsep}{4pt}
\begin{tabularx}{\linewidth}{@{}Xrrrc@{}}
\toprule
Проверка & Состояний & Глубина & Время, с & Итог \\
\midrule
\code{structure} & 14\,552\,302 & 110\,960 & 43,1 & PASS \\
\code{p01_crc_ignored} & 14\,552\,302 & 194\,839 & 62,3 & PASS \\
\code{p02_absent_slave} & 14\,552\,302 & 194\,839 & 69,3 & PASS \\
\code{p03_broadcast_silent} & 14\,552\,302 & 194\,839 & 74,7 & PASS \\
\code{p04_broadcast_written} & 14\,552\,302 & 194\,839 & 74,0 & PASS \\
\code{p05_addressed_write} & 14\,552\,302 & 194\,839 & 64,7 & PASS \\
\code{p06_write_isolation} & 14\,552\,302 & 194\,839 & 70,4 & PASS \\
\code{p07_normal_response} & 14\,552\,302 & 194\,839 & 67,0 & PASS \\
\code{p08_illegal_function} & 14\,552\,302 & 194\,839 & 72,4 & PASS \\
\code{p09_illegal_register} & 14\,552\,302 & 194\,839 & 66,9 & PASS \\
\code{p10_write_authorized} & 14\,552\,302 & 194\,839 & 69,2 & PASS \\
\code{p11_single_response} & 14\,552\,302 & 194\,839 & 68,7 & PASS \\
\code{p12_transaction_progress}\textsuperscript{*} & 28\,849\,450 & 194\,839 & 214,0 & PASS \\
\code{p13_completion_clean} & 14\,552\,302 & 194\,839 & 63,4 & PASS \\
\code{p14_invalid_broadcast} & 14\,552\,302 & 194\,839 & 62,0 & PASS \\
\bottomrule
\end{tabularx}}
\par\smallskip
{\footnotesize * Свойство p12 проверено отдельным запуском без ограничения времени Python-обертки.}
\end{table}


В таблице указано число сохраненных состояний из поля \code{states, stored}
и время, сообщаемое самим SPIN. При LTL-проверке исследуется поведение
модели совместно с автоматом свойства, поэтому число состояний может
различаться между формулами. Достигнутая глубина оказалась меньше
заданного предела в 1\,000\,000 шагов.

Большинство LTL-проверок потребовало от 62 до 74,7 секунды и около
1\,208 Мбайт памяти. Для свойства прогресса SPIN сообщил суммарное
использование 2\,358,511 Мбайт памяти. Общее время 15 успешных поисков
составило примерно 19 минут, без учета генерации, компиляции
и прерванной попытки проверки \code{p12}.

Полные логи находятся в папке запуска. Вместе с отчетом сохранены исходный
итог скрипта \code{verification/automated_report.txt}, лог отдельной
проверки \code{verification/p12_manual.log} и сводная статистика
\code{verification/summary.json}. Сводка успешных поисков учитывает
отдельную проверку \code{p12} и не заменяет исходный итог скрипта.

\subsubsection{Существенные уточнения модели и свойств}

Для проверки представленной версии изменять алгоритм Master и Slave
или переписывать LTL-формулы не потребовалось. Потребовалось уточнить
порядок выполнения верификации: свойство прогресса было проверено
отдельно, поскольку фиксированного времени Python-скрипта не хватило
для полного обхода.

При этом сама возможность содержательной проверки опирается на детали,
уже предусмотренные в модели. Снимки памяти позволяют сравнивать
состояния до и после запроса. Флаги записи отличают отсутствие операции
от повторной записи того же значения. Счетчик ответов сохраняет сведения
об отправке после извлечения сообщения из канала.

Три служебных уведомления задают момент, когда результат транзакции
действительно можно оценивать. Без ожидания завершения обработчиков
отсутствие ответа можно было бы принять за тайм-аут до того, как
адресованный Slave выполнил свой шаг. Состояние подготовки и условия
\code{active} и \code{finished} также отделяют наблюдения текущего
обмена от сброса признаков перед следующим запросом.

Следовательно, найденных противоречий между алгоритмом и сформулированными
требованиями нет. Результат относится к выбранным диапазонам данных,
начальному состоянию и последовательному выполнению транзакций.
Отсутствие ошибок в этой модели подтверждает описанную логику обработки,
но само по себе не является доказательством корректности полного
устройства Modbus.
